\section{Experiment Details}
\label{app:experiment}

\subsection{Data Construction}

\paragraph{Restaurant Selection.}
We select restaurants from the Yelp Open Dataset using a multi-stage pipeline:
(1) Filter by metropolitan area and target category (cafes);
(2) Rank candidates by review volume to ensure sufficient evidence;
(3) Apply LLM-based category verification to filter false positives.
We retain 20 restaurants meeting quality thresholds.

\paragraph{Review Sampling.}
For each restaurant, we construct a document set of up to 20 reviews.
To capture diverse user sentiments, we stratify sampling by star rating, targeting 4 reviews from each level (1--5 stars) where available.
Within each stratum, we select reviews by descending text length to maximize information density.
All reviewer identifiers are anonymized.

\subsection{Evidence Types}

We support four evidence types for evaluating user request conditions:

\paragraph{item\_meta.}
Direct lookup of structured item attributes (e.g., \texttt{attributes.WiFi}, \texttt{attributes.NoiseLevel}).
Missing values yield ``unknown'' (0); boolean attributes map directly to $\pm 1$.
Categorical attributes use declarative matching rules: if \texttt{true} condition matches, return $+1$; if \texttt{false} condition matches, return $-1$; otherwise return $0$.

\paragraph{item\_meta\_hours.}
Time-window containment verification.
Given business hours (e.g., ``7:00--18:00'') and required range (e.g., ``10:00--15:00''), returns $+1$ if the required range falls entirely within business hours, $-1$ otherwise.

\paragraph{review\_text.}
LLM-based review analysis.
An LLM judges each review independently for the target aspect, returning $+1$ (positive mention), $0$ (neutral/not mentioned), or $-1$ (negative mention).
Judgments are cached for reproducibility.
The final result uses \emph{symmetric consensus voting} with a 75\% threshold:
\begin{equation}
    \text{result} = \begin{cases}
        +1 & \text{if } \frac{n_+}{n_+ + n_-} \geq 0.75 \\
        -1 & \text{if } \frac{n_-}{n_+ + n_-} \geq 0.75 \\
        0  & \text{otherwise (no consensus)}
    \end{cases}
\end{equation}

\paragraph{review\_meta.}
Reviewer and review-level metadata conditions.
Supports filtering (e.g., elite reviewers only, reviews from 2018+, minimum star rating) and weighted aggregation (by review count, usefulness votes, or social connections).

\subsection{Request Structure}

Each request is represented as a Boolean expression tree.
Leaf nodes specify individual conditions with aspect names and evidence specifications.
Internal nodes apply Boolean operators (AND/OR) with three-valued semantics:

\begin{center}
\begin{tabular}{c|ccc}
AND & $+1$ & $0$ & $-1$ \\
\hline
$+1$ & $+1$ & $0$ & $-1$ \\
$0$  & $0$  & $0$ & $-1$ \\
$-1$ & $-1$ & $-1$ & $-1$
\end{tabular}
\quad
\begin{tabular}{c|ccc}
OR & $+1$ & $0$ & $-1$ \\
\hline
$+1$ & $+1$ & $+1$ & $+1$ \\
$0$  & $+1$ & $0$  & $0$ \\
$-1$ & $+1$ & $0$  & $-1$
\end{tabular}
\end{center}

\subsection{Request Groups}

\paragraph{G01 (R00--R04): Flat / Simple Metadata.}
Five requests using only direct item attributes with flat AND composition.
Example: ``Looking for a quiet cafe with free Wi-Fi and no TV.''

\paragraph{G02 (R05--R09): Mixed / Item Meta + Review Text.}
Five requests combining structured metadata filters with unstructured review text analysis.
Example: ``Budget-friendly indoor cafe where reviews mention good matcha.''

\paragraph{G03 (R10--R14): Flat / Computed Metadata.}
Five requests requiring hours interpretation.
Example: ``Quiet cafe with free Wi-Fi open Monday 10am--3pm.''

\paragraph{G04 (R15--R19): Flat / Social Metadata.}
Five requests using reviewer credibility signals.
Example: ``Places where experienced Yelp reviewers gave high ratings.''

\subsection{Ground Truth and Evaluation}

\paragraph{Label Generation.}
Ground truth labels are precomputed deterministically.
For each (restaurant, request) pair:
(1) Evaluate each condition against available evidence;
(2) Aggregate results using the request's Boolean structure;
(3) Store the final label ($+1$, $0$, or $-1$) with full evidence breakdown.
LLM judgments for review text are cached to ensure reproducibility across runs.
The resulting dataset contains 400 labeled pairs (20 restaurants $\times$ 20 requests).

\paragraph{Top-1 Selection.}
Each request is designed to have exactly one restaurant with a $+1$ (recommend) ground truth label.
This enables unambiguous evaluation: for each request, there is a single correct answer among the 20 candidates.

\paragraph{Hits@k Metric.}
We evaluate recommendation accuracy using Hits@k:
\begin{equation}
    \text{Hits@k} = \frac{1}{|R|} \sum_{r \in R} \mathbf{1}[\text{rank}(r) \leq k]
\end{equation}
where $R$ is the set of requests and $\text{rank}(r)$ is the position of the ground-truth positive item in the method's ranking for request $r$.
Hits@1 measures whether the method's top recommendation exactly matches the gold item.

\subsection{Baselines}
\label{app:baselines}

We evaluate against 13 zero-shot prompting baselines organized into two categories.

\subsubsection{Reasoning-Based Methods}

\paragraph{Chain-of-Thought (CoT).}
Zero-shot prompting where the LLM analyzes all candidates in a single context window and selects the best match.
Given a user request and the set of 20 candidate restaurants (each with metadata and reviews), the LLM is prompted to:
(1) analyze each restaurant's suitability for the user's stated requirements;
(2) select the single best-matching restaurant from the candidates.
No few-shot examples are provided, testing the model's inherent reasoning capabilities.

\paragraph{Self-Consistency (CoT-SC).}
Extends Chain-of-Thought by sampling multiple independent reasoning paths (typically 5) and selecting the final answer via majority vote.
This reduces sensitivity to any single reasoning chain's errors.

\paragraph{ReAct.}
Interleaves reasoning (Thought) with actions (examine reviews, check ratings, analyze requirements) in a multi-turn loop.
The model alternates between thinking about the problem and taking actions to gather information, with a maximum of 5 steps before making the final selection.

\paragraph{Least-to-Most (L2M).}
Decomposes the user request into a sequence of simpler subproblems, solves each subproblem incrementally (with solutions to earlier subproblems available as context), then synthesizes the final recommendation from all partial solutions.

\paragraph{Plan-and-Solve (PS).}
Generates an explicit multi-step plan before execution.
The model first devises a strategy for evaluating candidates, then systematically follows the plan to analyze each restaurant against the user's requirements.

\paragraph{Self-Ask.}
Generates follow-up questions about the user request (e.g., ``What specific attributes should I check for this requirement?''), answers each question using the available evidence, then uses the accumulated answers to make the final selection.

\paragraph{Decomposed Prompting.}
Uses specialized sub-task handlers orchestrated by a decomposer.
Handlers include: extract\_requirements (parse user constraints), analyze\_reviews (assess review sentiment), count\_evidence (tally supporting/opposing mentions), and make\_decision (synthesize final ranking).
Each handler focuses on a specific capability, improving modularity.

\subsubsection{Ranking-Based Methods}

\paragraph{Pointwise (Fine-Grained).}
Scores each candidate independently on a 4-point relevance scale: Highly Relevant (3), Relevant (2), Somewhat Relevant (1), Not Relevant (0).
Candidates are then ranked by their scores.
Complexity: $O(n)$ LLM calls for $n$ candidates.

\paragraph{PaRaDe.}
Few-shot demonstration-based ranking that provides explicit examples of how to compare and rank candidates.
The demonstrations guide the model toward consistent ranking behavior.

\paragraph{Pairwise (PRP).}
Compares all pairs of candidates with the prompt ``Which is more relevant to the user's request: A or B?''
Aggregates pairwise wins into a total score for each candidate, then ranks by win count.
Complexity: $O(n^2)$ LLM calls for $n$ candidates.

\paragraph{RankGPT.}
Sliding-window listwise reranking that processes candidates in overlapping windows.
The model ranks items within each window, and the process iterates to refine the global ranking.
Balances context length constraints with ranking quality.

\paragraph{Setwise.}
Compares candidates in sets (typically 3--5 items) rather than pairs.
For each set, the model identifies the best candidate; this process repeats with tournament-style elimination.
Offers a middle ground between pairwise granularity and listwise efficiency.

\paragraph{Listwise.}
Single-pass ranking where the LLM outputs a complete permutation of all candidates in one response.
The model sees all items simultaneously and produces the full ranking directly.
Complexity: $O(1)$ LLM calls (single call), but requires sufficient context length.
